\documentclass[11pt, a4paper]{article}
\usepackage[a4paper, margin=2.5cm]{geometry}
\usepackage{amsmath}
\usepackage{booktabs}
\usepackage{titlesec}
\usepackage{enumitem}
\usepackage{abstract}
\usepackage{authblk}
\usepackage[colorlinks=true, linkcolor=blue, citecolor=blue, urlcolor=blue]{hyperref}
\titleformat{\section}{\large\bfseries}{\thesection}{1em}{}
\titleformat{\subsection}{\normalsize\bfseries}{\thesubsection}{1em}{}
\title{\textbf{The BX3 Framework: A Universal Architecture of Functional Roles for Purpose, Bounded Reasoning, and Deterministic Fact}}
\author{Jeremy Blaine Thompson Beebe}
\affil{Independent Researcher}
\date{April 2026}
\begin{document}
\maketitle
%====================================================================
% EDITOR'S NOTE — v2 CORRECTIONS
%====================================================================
% 1. Removed false field-trial claim (was June 2026, impossible post-April submission)
% 2. Replaced with: Agentic runtime (in development) + Irrig8 product use case
% 3. Added TAO distinction, governance-translation citation, Sandbox Gate clarification
% 4. Fixed April/June 2026 timeline contradiction
%====================================================================
\section*{Abstract}
The current discourse around artificial intelligence frequently presents a false binary: either AI will replace traditional software systems, or it is merely a productivity tool of limited consequence. This paper argues that both positions miss a more fundamental and practically useful insight. We propose the BX3 Framework --- a universal architectural model organized around three immutable functional layers: the Purpose Layer, which provides intent, judgment, and accountability; the Bounds Engine, which provides interpretation, bounded reasoning, and constrained execution; and the Fact Layer, which provides deterministic enforcement, hard physical constraint, and forensic auditability. Critically, the BX3 Framework does not prescribe who or what occupies each layer. It prescribes the functional properties each layer must maintain --- and holds any actor occupying that layer to those properties regardless of their nature. A human, an AI system, or a mechanical process may each legitimately occupy any layer, provided they satisfy that layer's functional requirements. This actor-agnostic, function-centered definition makes the framework universally applicable across human organizations, fully automated systems, multi-agent AI architectures, and any hybrid composition. The BX3 Framework is currently being implemented as the Agentic runtime --- an open reference architecture for building BX3-compliant systems --- with a first production application in precision agriculture (Irrig8, an autonomous farming OS for resource health management and forensic accountability). Field trials are planned upon completion of the Agentic v1.0 runtime. A companion engineering specification --- the BX3 Protocol --- is in preparation and will provide normative, certifiable implementation standards derived from this theoretical foundation.
\section{Introduction}
The current discourse around artificial intelligence frequently presents a false binary: either AI will replace traditional software systems, or it is merely a productivity tool of limited consequence. This paper argues that both positions miss a more fundamental and practically useful insight. We propose the BX3 Framework --- a universal architectural model organized around three immutable functional layers: the Purpose Layer, which provides intent, judgment, and accountability; the Bounds Engine, which provides interpretation, bounded reasoning, and constrained execution; and the Fact Layer, which provides deterministic enforcement, hard physical constraint, and forensic auditability. Critically, the BX3 Framework does not prescribe who or what occupies each layer. It prescribes the functional properties each layer must maintain --- and holds any actor occupying that layer to those properties regardless of their nature. A human, an AI system, or a mechanical process may each legitimately occupy any layer, provided they satisfy that layer's functional requirements. This actor-agnostic, function-centered definition makes the framework universally applicable across human organizations, fully automated systems, multi-agent AI architectures, and any hybrid composition. A core safety property of the framework is its upstream accountability guarantee: when any node encounters a state it cannot resolve within its bounds, accountability escalates recursively upward through the system hierarchy --- bypassing all machine actors --- until it reaches a human anchor. The system fails upward into human consciousness --- never downward into algorithmic chaos. We demonstrate that when these three functional layers are clearly separated and their properties enforced by design, systems become simultaneously more capable and less complex. The BX3 Framework is being implemented as the Agentic runtime architecture, with Irrig8 --- an end-to-end autonomous farming OS for resource health management and auditability, targeting center-pivot irrigation in Colorado's San Luis Valley --- as the first production product use case. Field trials are planned following the v1.0 release of the Agentic runtime. A companion engineering specification --- the BX3 Protocol --- is in preparation and will provide normative, certifiable implementation standards derived from this theoretical foundation.
\section{The Three Immutable Functional Layers}
\subsection{Purpose Layer}
The Purpose Layer provides intent, judgment, and accountability. It is the layer at which goals are set, trade-offs are weighed, and final accountability is assigned. The Purpose Layer does not execute actions; it defines what ought to happen and is responsible for what does happen. Any actor occupying the Purpose Layer must be capable of bearing final accountability. In human organizations, this layer is occupied by boards, executives, regulators, and civil institutions. In autonomous systems, this layer must be occupied by an actor capable of bearing genuine accountability --- a requirement that, for all current and foreseeable AI systems, means a human or a human institution. The key functional property of the Purpose Layer is \textit{accountability without plausible deniability}. An actor in this layer cannot claim it was merely executing instructions when those instructions produced harm.
\subsection{Bounds Engine}
The Bounds Engine provides interpretation, bounded reasoning, and constrained execution. It is the layer that analyzes context, proposes actions, and exercises judgment within constraints. The Bounds Engine is intentionally limbless --- it has no direct actuators, no physical effectors, and no ability to act on the world outside the interpretive and analytical domain. It may produce analysis, recommendations, and constrained decisions, but any action it proposes must be validated by the Fact Layer before execution. The key functional property of the Bounds Engine is \textit{bounded agency}: it can reason and recommend but cannot independently execute. When the Bounds Engine encounters a state it cannot resolve within its bounds, it must escalate to the Purpose Layer --- not attempt to expand its own authority.
\subsection{Fact Layer}
The Fact Layer provides deterministic enforcement, hard physical constraint, and forensic auditability. It is the layer that interacts with the physical world through actuators, sensors, and verified state transitions. The Fact Layer is the only layer that can produce non-reversible physical effects. Its defining characteristic is 100\% deterministic behavior: given the same inputs, it must produce identical outputs, and it cannot vary its behavior based on probability, context, or learned weights. The Fact Layer enforces constraints that are not merely business rules or software permissions but physical limits --- circuit breakers, mechanical interlocks, chemical process controls, and access restrictions that exist independent of any computational system. The Fact Layer maintains a forensic ledger of all state changes, providing a complete and tamper-evident record of all physical actions.
\section{The Five Architectural Pillars}
\subsection{Loop Isolation}
The Bounds Engine and the Fact Layer never share a functional plane. This means that a single AI agent cannot simultaneously propose and execute an action. The Bounds Engine may analyze, recommend, and decide within its bounded domain, but any action that crosses into the physical world must be mediated by the Fact Layer, which operates under deterministic and formally verifiable constraints.
\subsection{Recursive Spawning and Worksheet Architecture}
Systems built on the BX3 Framework support recursive spawning: a parent node can birth a child loop that carries a hard-coded pointer to the parent's Purpose Layer. Each child loop operates within a containerized ``Worksheet'' that preserves local state and local goal structures while remaining subordinate to the parent's accountability. This enables the system to maintain goal coherence across arbitrarily long chains of delegated reasoning and action without requiring continuous cloud connectivity or centralized control.
\subsection{Spatial Firewall}
The access tiers of a BX3-compliant system are not implemented as software permissions but as physical Fact Layer constraints. A user who lacks authorization cannot be granted access by overriding software alone --- the physical access infrastructure must be reconfigured. This distinguishes BX3 spatial controls from conventional RBAC and ACL models.
\subsection{Root Tunneling}
A human operator can project their Purpose Layer authority into any child node or Worksheet at any depth, effectively ``tunneling'' through the system hierarchy to assume direct accountability for a specific decision or action. This ensures that the accountability guarantee is not merely a structural property but an operational capability.
\subsection{Bailout Protocol}
When any Bounds Engine encounters a state it cannot resolve within its current authority --- a novel situation, a conflict between competing imperatives, or an action that would exceed its operational constraints --- it invokes the Bailout Protocol. The protocol immediately suspends autonomous action and escalates the decision to the nearest human with appropriate authority, bypassing all intermediate machine actors. This is architecturally distinct from Tiered Agentic Oversight (TAO) models, which escalate errors between successive AI tiers. In TAO, each tier attempts to resolve the error before passing it up; in BX3's Bailout Protocol, the human anchor is contacted \textit{directly} with full diagnostic context, with all machine intermediaries relegated to relay status. The Bailout Protocol also hard-blocks any Bounds Engine proposal that would violate pre-defined safety or physical constraints.
\section{Governance and Compliance}
The BX3 Framework is designed to be compatible with and supportive of existing governance standards, not in competition with them. The Purpose Layer maps naturally onto organizational accountability structures, regulatory compliance frameworks, and legal personhood requirements. The Fact Layer satisfies deterministic audit requirements that regulators in aviation, medical devices, financial systems, and safety-critical infrastructure already apply to their certified systems. The Bounds Engine provides the interpretive and adaptive capability that allows systems to operate effectively in novel situations while remaining within the bounds set by the Purpose Layer. This architectural claim is consistent with recent work on translating governance norms into enforceable controls. The distinction between BX3's Sandbox Gate and conventional sandboxing should be clarified: BX3's Sandbox Gate requires a mandatory digital twin simulation with explicit human approval before physical actuators are unlocked, where generic sandboxing refers to isolated execution environments without mandatory physical-world confirmation requirements.
\section{Formal Properties and Proof Sketch}
We sketch the formal properties that a BX3-compliant system must satisfy. The key invariants are: (1) \textbf{Accountability Invariant}: Every action in the Fact Layer is traceable to a Purpose Layer actor. (2) \textbf{Bounds Invariant}: The Bounds Engine never executes an action that has not been validated by the Fact Layer. (3) \textbf{Escalation Invariant}: Any state that the Bounds Engine cannot resolve within its current authority triggers the Bailout Protocol, reaching a human anchor without passing through intermediate machine actors. (4) \textbf{Determinism Invariant}: The Fact Layer produces identical outputs for identical inputs, enabling formal verification and forensic audit.
\section{Determinism, Hybrid Systems, and the Infrastructure Argument}
\subsection{The Value of Determinism is Not Compensable}
Determinism provides properties that probabilistic systems cannot replicate: reproducibility (enabling debugging, auditing, and legal accountability), formal verification (mathematical proof of properties under all inputs), certification (regulatory frameworks in aviation, medical devices, financial systems, and safety-critical infrastructure require deterministic behavior), and latency (hard real-time requirements are achievable with deterministic systems and not with current AI inference pipelines).
\subsection{The Infrastructure Argument}
Every AI system in production today runs on top of vast deterministic infrastructure: operating systems, databases, networking stacks, authentication systems, billing pipelines, and monitoring tools. The claim that AI will replace software is structurally self-refuting --- the AI systems making this possible are themselves dependent on deterministic software that will not and should not be replaced.
\subsection{The Regulatory and Trust Barrier}
Even in domains where AI could theoretically replace deterministic systems on a performance basis, the practical barriers are substantial. Organizations with core operations dependent on software will not replace functioning, auditable, certified systems with probabilistic alternatives without extraordinary evidence of equivalent reliability and auditability.
\section{Relationship to Prior Work}
The BX3 Framework synthesizes several established frameworks rather than claiming wholly original invention. Separation of Concerns (Dijkstra): Extended here to explicitly include the separation of actors based on their functional roles rather than their physical or computational nature. Zero Trust Architecture: BX3's Spatial Firewall can be viewed as a physical-world extension of Zero Trust principles, where the ``never trust, always verify'' maxim is applied to physical access rather than network access. Tiered Agentic Oversight (TAO): BX3's Bailout Protocol is architecturally distinct from TAO. In TAO, errors escalate between successive AI tiers; in BX3, the Bailout Protocol bypasses all machine actors and contacts the human anchor directly with full diagnostic context. Governance Norms to Enforceable Controls: BX3's architectural claim is consistent with recent work translating governance norms into enforceable controls, as both seek to operationalize accountability at the system level.
\section{Conclusion}
The BX3 Framework offers a path beyond the false binary of AI replacement versus AI as limited tool. By defining three immutable functional layers with enforced properties, and by guaranteeing that accountability always reaches a human, the framework enables systems that are simultaneously more capable and more controllable. The framework is currently being implemented as the Agentic runtime, with Irrig8 as the first production product use case.
\begin{thebibliography}{99}
\item E.W. Dijkstra, ``On the role of scientific thought,'' 1982.
\item NIST, \textit{Artificial Intelligence Risk Management Framework (AI RMF 1.0)}, 2023.
\item ISO/IEC 42001:2023, \textit{Information Technology --- Artificial Intelligence}, 2023.
\item Google DeepMind, \textit{Tiered Agentic Oversight (TAO)}, arXiv:2506.12482, June 2025.
\item \textit{Governance Norms to Enforceable Controls}, arXiv:2604.05229, April 2026.
\item \textit{Type-Checked Compliance / Lean-Agent Protocol}, arXiv:2604.01483, April 2026.
\item \textit{Agentic AI Evolution}, arXiv:2602.10479, February 2026.
\item \textit{Three-Layer Global AI Governance}, Lawfare, February 2026.
\end{thebibliography}
\end{document}
